\begin{lemma}{Eigenschaften des k--ten Maximums}
             {EigenschaftenDesKtenMaximums}
Sei $n \in \N$.\\
Dann gilt:
\begin{enumerate}[(1)]
	\item \label{EigenschaftenDesKtenMaximums1}
	$\forall\, k \in \haken {n-1}: \mu_{k+1,n} \le \mu_{k,n}$.
	\item \label{EigenschaftenDesKtenMaximums2}
	$\forall\, k \in \haken n\, \forall\, x \in \R^n\, \forall\,a \in \R:
	\mu_{k,n}(x) \ge a \iff \abs{\menge{i \in \haken n}{x_i \ge a}} \ge k$.
	\item\label{EigenschaftenDesKtenMaximums2a}
	$\forall\, k \in \haken n\, \forall\, x \in \R^n\, \forall\,a \in \R:
	\mu_{k,n}(x) < a \iff \abs{\menge{i \in \haken n}{x_i < a}} > n - k$.
	\item\label{EigenschaftenDesKtenMaximums2b}
	$\forall\, k \in \haken n\, \forall\, x \in \R^n\, \forall\,a \in \R:
	\mu_{k,n}(x) > a \iff \abs{\menge{i \in \haken n}{x_i > a}} \ge k$.
	\item\label{EigenschaftenDesKtenMaximums2c}
	$\forall\, k \in \haken n: \mu_{k, n}$ stetig.
	\item \label{EigenschaftenDesKtenMaximums3}
	$\forall\, k \in \haken n: \mu_{k, n}$ nichtnegativ--homogen.
\end{enumerate}
\end{lemma}
\begin{proof}
(1) Seien $k \in \haken {n-1}$ und $x \in \R^n$.\\
Seien
\[M_{k, x}\coloneqq \menge{x_i}{i\in\haken n,\abs{\menge{j\in\haken n}{x_j>x_i}}<k}\]
und
\[M_{k+1,x}\coloneqq
\menge{x_i}{i\in\haken n,\abs{\menge{j\in\haken n}{x_j>x_i}}<k+1}.\]
Nach~\refclaim{Sortierungsmengen}{Sortierungsmengen1} ist dann
$M_{k,x} \subseteq M_{k+1,x}$ und daher:
\[\mu_{k+1,n}(x) = \min M_{k+1,x} \le \min M_{k,x} = \mu_{k,n}(x).\]
(2) Seien $k \in \haken n$, $x \in \R^n$ und $a \in \R$.\\
\enquote{$\Rightarrow$}: Es gelte $\mu_{k,n}(x) \ge a$.\\
Nach (1) gilt: $\forall\,j\in\haken k: \mu_{j, n}(x) \ge \mu_{k, n}(x)\ge a$.\\
Es gibt $\iota: \haken k \rightarrow \haken n$ injektiv mit
$\forall\, j \in \haken k: x_{\iota(j)}= \mu_{j,n}(x)$.\\
Ferner gilt $\iota(\haken k) \subseteq \menge{i \in \haken n}{x_i \ge a}$ und
damit wegen der Injektivität von $\iota$
\[k = \abs{\iota(\haken k)} \le \abs{\menge{i \in \haken n}{x_i \ge a} }.\]
\enquote{$\Leftarrow$:} Es gelte $\mu_{k,n}(x) < a$.\\
Es gibt ein $j \in \haken n$ mit $x_j = \mu_{k,n}(x) < a$.\\
Nach Definition von $\mu_{k,n}$ gilt weiter: 
$\abs{\menge{i \in \haken n}{x_i> x_j}}<k$.\\
Somit gilt 
$\menge{i \in \haken n}{x_i\ge a} \subseteq \menge{i \in \haken n}{x_i > x_j}$
und daher
\[\abs{\menge{i\in\haken n}{x_i\ge a}}
\le\abs{\menge{i\in\haken n}{x_i>x_j}}<k.\]
(3) Seien $k \in \haken n$, $x \in \R^n$ und $a \in \R$. Dann gilt:
\begin{align*}
\mu_{k, n}(x) < a & \overset{(2)}\iff \abs{\menge{i \in \haken n}{x_i \ge a}}<k
\text.
\end{align*}
Offensichtlich gilt auch:
\[\abs{\menge{i \in \haken n}{x_i \ge a}}<k \iff
  \abs{\menge{i \in \haken n}{x_i < a}} > n-k\text.\]
Damit ist (3) gezeigt.\\
(4) Analog zu (2).\\
(5) Es genügt zu zeigen, dass Urbilder offener Intervalle offen sind.\\
Seien also $a, b \in \R$ mit $a < b$ und $k \in \haken n$. Dann gilt:
\begin{align*}
&\menge{x \in \R^n}{\mu_{k, n}(x)> a}
\overset{(4)}= \menge{x \in \R^n}{\exists A \subseteq \haken n: \abs A
\ge k \wedge (\forall\, i \in A: x_i > a)}\\
=\,&\bigcup_{\substack{A \subseteq \haken n,\\ \abs A \ge k}}
\bigcap_{i \in \haken A}\menge{x \in \R^n}{x_i> a}
= \bigcup_{\substack{A \subseteq \haken n,\\ \abs A \ge k}}
\bigcap_{i \in \haken A} \pr_i\inv(]a, \infty[) \eqqcolon G\text.
\end{align*}
Da $\pr_i$ für alle $i \in \haken n$ unabhängig von der Norm stetig sind, ist
$G$ als Vereinigung endlicher Schnitte offener Mengen offen.\\
Analog gilt:
\begin{align*}
&\menge{x \in \R^n}{\mu_{k, n}(x)< b}
\overset{(3)}= \menge{x \in \R^n}{\exists B \subseteq \haken n: \abs B
>n- k \wedge (\forall i \in B: x_i < b)}\\
=&\bigcup_{\substack{B \subseteq \haken n,\\ \abs B > n - k}}
\bigcap_{i \in \haken B}\menge{x \in \R^n}{x_i< b}
= \bigcup_{\substack{B \subseteq \haken n,\\ \abs B > n - k}}
\bigcap_{i \in \haken B} \pr_i\inv(]-\infty, b[) \eqqcolon K\text.
\end{align*}
Auch $K$ ist (mit dem gleichen Argument wie $G$) offen.
Außerdem ist
\[\mu_{k,n}\inv(]a,b[) = \menge{x \in \R^n}{\mu_{k,n}(x) > a \wedge
\mu_{k,n}(x) < b} = G \cap K\]
und damit offen, weil $G$ und $K$ es sind.\\
(6) Seien $k \in \haken n$, $x \in \R^n$ und $t \in \R_{>0}$.\\
Dann gilt:
\begin{align*}
&\menge{(tx)_i}{i\in\haken n,\abs{\menge{j\in\haken n}{(tx)_j>(tx)_i}}<k}\\
=&\,\menge{tx_i}{i\in\haken n,\abs{\menge{j\in\haken n}{tx_j>tx_i}}<k}\\
\overset{t>0}= &\,
\menge{tx_i}{i\in\haken n,\abs{\menge{j\in\haken n}{x_j>x_i}}<k}\\
=&\, t\menge{x_i}{i\in\haken n,\abs{\menge{j\in\haken n}{x_j>x_i}}<k}\text.
\end{align*}
Also nach Definition:
\begin{align*}
\mu_{k, n}(tx) &=
\min\menge{(tx)_i}{i\in\haken n,\abs{\menge{j\in\haken n}{(tx)_j>(tx)_i}}<k}\\
&=\min
\left(t\menge{x_i}{i\in\haken n,\abs{\menge{j\in\haken n}{x_j>x_i}}<k}\right)
\overset{t>0}=t\mu_{k, n}(x)\text.
\end{align*}
Zudem gilt offenbar $\mu_{k,n}(0) = 0$, sodass $\mu_{k,n}$
nichtnegativ--homogen ist.
\end{proof}